\section{Experimental Setup}
\label{sec:experiments}

Before presenting the empirical comparison, we summarize the experimental protocols used to evaluate the one-step methods surveyed in this paper. Because the field is evolving rapidly, exact hyperparameters and architectural details differ across works; our aim is to provide a unified description of the settings in which the results of \Cref{sec:results} were obtained.

% ---------------------------------------------------------------------
\subsection{Datasets and Metrics}

The de facto benchmark for one-step generative modeling is class-conditional image generation on ImageNet \citep{deng2009imagenet}. Most recent methods report results at both $256\times256$ and $512\times512$ resolutions, training on the standard ImageNet training split and evaluating on the validation split. The primary evaluation metric is the Fréchet Inception Distance (FID) \citep{heusel2017gans}, which measures the distributional distance between generated and real image features extracted by an Inception-V3 network \citep{szegedy2016rethinking}. Lower FID indicates better sample quality.

Additional metrics reported in the literature include Inception Score (IS) \citep{salimans2016improved}, precision and recall \citep{kynkaanniemi2019improved}, and the number of function evaluations (NFE) required per sample. For flow-based methods, NFE is determined by the number of neural-network evaluations during numerical integration; for one-step methods, NFE is ideally one.

% TODO(team): Add specific dataset details (train/val sizes, preprocessing) and any non-image benchmarks if applicable.

% ---------------------------------------------------------------------
\subsection{Training and Evaluation Protocol}

Following the experimental protocol of \citet{zhang2025alphaflow}, we train all three velocity-based methods from scratch on class-conditional ImageNet-1K 256$\times$256. We use the vanilla DiT architecture at three scales: DiT-B/2 (131M parameters), DiT-XL/2 (676M parameters), and a fine-tuned DiT-XL/2+ variant. The B/2 and XL/2 models are trained for 240 epochs; XL/2+ is initialized from the XL/2 checkpoint and fine-tuned for an additional 60 epochs with a larger batch size of 1024. All models operate in the latent space of the Stable Diffusion VAE.

The optimizer is AdamW with a learning rate schedule tuned for each method. MeanFlow and AlphaFlow rely on Jacobian--vector product (JVP) computations to estimate the total derivative of the average velocity; AlphaFlow additionally implements a sigmoid curriculum over the consistency step ratio $\alpha$, annealing from trajectory flow matching ($\alpha=1$) toward MeanFlow ($\alpha\to 0$) with a clamping value of $\eta=5\times 10^{-3}$. Consistency-FM is trained with its velocity-consistency objective; because we have not yet completed Consistency-FM runs, its results in \Cref{tab:imagenet-results} are reported as placeholders.

Evaluation follows \citet{zhang2025alphaflow}: we generate 50,000 class-conditional samples and report the Fr\'{e}chet Inception Distance (FID) \citep{heusel2017gans} and the Fr\'{e}chet DINOv2 Distance (FDD). We evaluate both one-step (NFE=1) and two-step (NFE=2) generation. For two-step sampling we use either ODE or consistency sampling, selecting the intermediate timestep that balances FID and FDD.

% ---------------------------------------------------------------------
\subsection{Baselines}

The main empirical comparison in \Cref{sec:results} contrasts the three surveyed velocity-based methods---Consistency-FM, MeanFlow, and AlphaFlow---trained at the DiT-B/2, DiT-XL/2, and DiT-XL/2+ scales. In addition, we report representative few-step baselines from the AlphaFlow paper that use the same DiT-XL/2 backbone and are trained from scratch, including Shortcut-XL/2, IMM-XL/2, MeanFlow-XL/2, MeanFlow-XL/2+, and FACM-XL/2. These baselines provide a direct parameter-matched comparison against our XL/2 and XL/2+ models.

% TODO(team): Add citation keys for Shortcut, IMM, and FACM if desired.

% ---------------------------------------------------------------------
\subsection{Computational Resources}

All models are trained on clusters of NVIDIA A100 or H100 GPUs. The dominant cost for MeanFlow and AlphaFlow comes from the JVP computation in the backward pass and the larger effective batch size used for XL/2+ fine-tuning. Consistency-FM avoids JVPs but requires EMA and consistency-based training. We report total training epochs and parameter counts in \Cref{tab:imagenet-results} to enable fair comparisons across methods.
